\subsection{Noether 环的同调维数}

设 A 为 Noether 局部环，M 为有限生成 A-模。回忆 (A, X, § 3, 第 6 节)，M 的一个分解

\[
\dots \longrightarrow \mathrm{L} _ {n} \stackrel {d _ {n}} {\longrightarrow} \mathrm{L} _ {n - 1} \longrightarrow \dots \longrightarrow \mathrm{L} _ {0} \stackrel {d _ {0}} {\longrightarrow} \mathrm{M} \rightarrow 0
\]

称为极小投射分解，如果每个模 \(L_{i}\) 都是有限型自由模，且复形 \(\kappa_{A} \otimes_{A} L\) 的微分为零。对每个整数 \(i \geqslant 0\)

我们有

\[
[ \mathrm{Ext} _ {\mathrm{A}} ^ {i} (\mathrm{M}, \kappa_ {\mathrm{A}}): \kappa_ {\mathrm{A}} ] = [ \mathrm{Tor} _ {i} ^ {\mathrm{A}} (\mathrm{M}, \kappa_ {\mathrm{A}}): \kappa_ {\Lambda} ] = \mathrm{rg} _ {\mathrm{A}} (\mathrm{L} _ {i})\tag{1}
\]

(A, X, 第 103 页, 例 3)。每个有限生成 A-模都具有这样一个分解 (A, X, 第 56 页, 命题 10)。

命题 4. 设 A 为 Noether 局部环，M 为有限生成 A-模，n 为一个 ≥ 0 的整数。下列条件等价：

\begin{enumerate}
\def\labelenumi{(\roman{enumi})}
\item
  有 \(\mathrm{dp}_{\mathrm{A}}(\mathrm{M}) < n\)
\item
  有 \(\mathrm{Tor}_n^{\mathrm{A}}(\mathrm{M},\kappa_{\mathrm{A}}) = 0\)；
\item
  有 \(\operatorname{Ext}_{\mathrm{A}}^{n}(\mathrm{M},\kappa_{\mathrm{A}}) = 0\)；
\item
  M 的每个极小投射分解的长度都 \textless{} n。
\end{enumerate}

蕴含关系 (i)⇒(ii) 与 (i)⇒(iii) 是直接的 (A, X, 第 100 页, 定理 1)。设 L 为 M 的一个极小投射分解；若 (ii) 或 (iii) 成立，则由 (1) 有 \(L_{n}=0\)。由于 M 的每个极小投射分解都同构于 L (A, X, 第 54 页, 命题 8)，我们得到 (iv)。蕴含关系 (iv)⇒(i) 是平凡的。

推论 1.--- 设 A 为 Noether 局部环，n 为一个 ≥ 0 的整数。下列条件等价：

\begin{enumerate}
\def\labelenumi{(\roman{enumi})}
\item
  有 \(\mathrm{dh}(\mathrm{A}) < n\)
\item
  有 \(\mathrm{Ext}_{A}^{i}(\mathrm{M},\mathrm{N}) = 0\) 且 \(\mathrm{Tor}_i^A (\mathrm{M},\mathrm{N}) = 0\)（对每对 A-模 \((\mathbf{M},\mathbf{N})\) 及每个整数 \(i\geqslant n\)）；
\item
  有 \(\mathrm{Tor}_n^{\mathrm{A}}(\kappa_{\mathrm{A}},\kappa_{\mathrm{A}}) = 0\)
\item
  有 \(\operatorname{Ext}_{\mathrm{A}}^{n}(\kappa_{\mathrm{A}},\kappa_{\mathrm{A}}) = 0\)；
\item
  有 \(\mathrm{dp}_{\mathrm{A}}(\kappa_{\mathrm{A}}) < n\)。
\end{enumerate}

显然 (i) 蕴含 (ii)，且 (ii) 蕴含 (iii) 与 (iv)。对 A-模 \(\kappa_{\mathrm{A}}\) 应用命题 4，可知 (iii) 与 (iv) 各自蕴含 (v)。我们证明 (v) 蕴含 (i)：若 \(\mathrm{dp}_{A}(\kappa_{\mathrm{A}}) < n\)，则对每个 A-模 M 有 \(\operatorname{Tor}_{n}^{\mathrm{A}}(\mathrm{M}, \kappa_{\mathrm{A}}) = 0\)；于是每个有限生成 A-模都具有有限投射维数 \(< n\) (命题 4)，这蕴含 \(\mathrm{dh}(\mathrm{A}) < n\) (A, X, 第 138 页, 命题 4)。

\[
\text { COROLLARY   2. } - \text { We   have   dh(A) = dp } _ {A} (\kappa_ {A})  .
\]

注.- 1) 设 A 为局部环。A-模 \(\mathrm{Tor}_{1}^{\mathrm{A}}(\kappa_{\mathrm{A}},\kappa_{\mathrm{A}})\) 同构于 \(\mathfrak{m}_{A}/\mathfrak{m}_{A}^{2}\)。因此，当 A 是 Noether 的时，为使 \(\mathrm{Tor}_{1}^{\mathrm{A}}(\kappa_{\mathrm{A}},\kappa_{\mathrm{A}})\) 为零，必须且只需 \(\mathfrak{m}_{\mathrm{A}}\) 为零，即 A 是一个域。于是推论 1 蕴含：同调维数为 0 的 Noether 局部环是域。

\begin{enumerate}
\def\labelenumi{\arabic{enumi})}
\setcounter{enumi}{1}
\tightlist
\item
  设 A 为 Noether 局部环，M 为具有有限投射维数 n 的有限生成 A-模，N 为非零的有限生成 A-模。A-模 Ext \(_{A}^{n}(M,N)\) 非零：事实上，设 L 为 M 的一个极小投射分解，d 为其微分。我们有正合列
\end{enumerate}

\[
\operatorname{Hom} _ {\mathrm{A}} \left(\mathrm{L} _ {n - 1}, \mathrm{N}\right) \xrightarrow {\operatorname{Hom} \left(d _ {n} , 1\right)} \operatorname{Hom} _ {\mathrm{A}} \left(\mathrm{L} _ {n}, \mathrm{N}\right) \longrightarrow \operatorname{Ext} _ {\mathrm{A}} ^ {n} (\mathrm{M}, \mathrm{N}) \rightarrow 0.
\]

由于 \(d_n \otimes 1_{\kappa_A}\) 为零，用 \(\kappa_A\) 作张量积得到一个同构 \(\kappa_A \otimes_A \operatorname{Hom}_A(L_n, N) \longrightarrow \kappa_A \otimes_A \operatorname{Ext}_A^n(M, N)\)，于是，考虑到上面的公式 (1)，

\[
\left[ \kappa_ {\mathrm{A}} \otimes_ {\mathrm{A}} \operatorname{Ext} _ {\mathrm{A}} ^ {n} (\mathrm{M}, \mathrm{N}): \kappa_ {\mathrm{A}} \right] = \left[ \operatorname{Ext} _ {\mathrm{A}} ^ {n} (\mathrm{M}, \kappa_ {\mathrm{A}}): \kappa_ {\mathrm{A}} \right] \left[ \kappa_ {\mathrm{A}} \otimes_ {\mathrm{A}} \mathrm{N}: \kappa_ {A} \right],
\]

由命题 4 与 Nakayama 引理，它非零。因此，M 的投射维数是使 Ext \(_{A}^{i}(M,N)\) 非零的最大整数 i。

\begin{enumerate}
\def\labelenumi{\arabic{enumi})}
\setcounter{enumi}{2}
\tightlist
\item
  设 A 为 Noether 环，M 为具有有限投射维数的有限生成 A-模，N 为支撑集等于 \(\operatorname{Spec}(A)\) 的有限生成 A-模。由前注及第 2 节的命题 2 与命题 3，投射维数 \(n\)（即 \(M\) 的）是这样的最大整数 \(i\)：\(\operatorname{Ext}_{A}^{i}(M,N)\) 非零；A-模 \(\operatorname{Ext}_{A}^{n}(M,N)\) 的支撑集是元素 \(\mathfrak{p}\)（属于 \(\operatorname{Spec}(A)\)）组成的集合，这里 \(\mathrm{dp}_{A_{\mathfrak{p}}}(M_{\mathfrak{p}})=n\)。
\end{enumerate}

可能存在这样的有限生成 A-模 M，其有限投射维数为 \(+\infty\)，且满足 \(\operatorname{Ext}_{\mathrm{A}}^{i}(\mathrm{M},\mathrm{A})=0\)（对充分大的 \(i\)）：*例如，A-模 \(\kappa_{\Lambda}\) 当 \(\Lambda\) 是非正则的 Gorenstein 局部环时即属此情形*。

命题 5.--- 设 A 为 Noether 环，M 为有限生成 A-模，n 为一个 ≥ 0 的整数。下列条件等价：

\begin{enumerate}
\def\labelenumi{(\roman{enumi})}
\item
  有 \(\mathrm{dp}_{\mathrm{A}}(\mathrm{M}) < n\)
\item
  对每个极大理想 \(\mathfrak{m}\)（属于 \(A\)），有 \(\operatorname{Ext}_{\mathrm{A}}^{n}(\mathrm{M},\mathrm{A}/\mathfrak{m}) = 0\)（相应地，有 \(\operatorname{Tor}_{n}^{\mathrm{A}}(\mathrm{M},\mathrm{A}/\mathfrak{m}) = 0\)）；
\item
  对 A 的每个极大理想 m，有 \(\operatorname{Ext}_{\mathrm{A}_{m}}^{n}(\mathrm{M}_{\mathfrak{m}},\mathrm{A}/\mathfrak{m})=0\)（相应地，有 \(\operatorname{Tor}_{n}^{\mathrm{A}_{\mathfrak{m}}}(\mathrm{M}_{\mathfrak{m}},A/\mathfrak{m})=0\)）。
\end{enumerate}

(i)⇒(ii)：这是显然的。

(ii)⇒(iii)：这由第 2 节的命题 2 得出。

\begin{enumerate}
\def\labelenumi{(\roman{enumi})}
\setcounter{enumi}{2}
\tightlist
\item
  (iii) \(\Rightarrow\) (i)：由命题 4，条件 (iii) 蕴含不等式 \(\mathrm{dp}_{A_{\mathfrak{m}}}(M_{\mathfrak{m}}) < n\) 对 A 的每个极大理想 \(\mathfrak{m}\) 成立，于是由命题 3 即得结论。
\end{enumerate}

注 4. — 设 A 为 Noether 环，\(n\) 为一个 \(\geqslant 0\) 的整数。若 \(\mathfrak{m}\) 与 \(\mathfrak{m}'\) 是 A 的两个不同的极大理想，则 A-模 \(\mathrm{Ext}_{\mathrm{A}}^{n}(\mathrm{A/m},\mathrm{A/m}')\) 与 \(\mathrm{Tor}_{n}^{\mathrm{A}}(\mathrm{A/m},\mathrm{A/m}')\) 都被 \(\mathfrak{m} + \mathfrak{m}'\) 零化，故均为零。用与命题 4 的推论 1 类似的证明，由命题 5 推出下列条件的等价性：

\begin{enumerate}
\def\labelenumi{(\roman{enumi})}
\item
  有 \(\mathrm{dh}(\mathrm{A}) < n\)
\item
  有 \(\operatorname{Ext}_A^i (\mathrm{M},\mathrm{N}) = 0\) 且 \(\operatorname{Tor}_i^{\mathrm{A}}(\mathrm{M},\mathrm{N}) = 0\)（对每对 A-模 \((\mathrm{M},\mathrm{N})\) 及每个整数 \(i\geqslant n\)）；
\item
  有 \(\mathrm{Tor}_n^{\mathrm{A}}(\mathrm{A} / \mathfrak{m},\mathrm{A} / \mathfrak{m}) = 0\)（对 A 的每个极大理想 \(\mathfrak{m}\)）；
\item
  有 \(\operatorname{Ext}_{\mathsf{A}}^{n}(\mathsf{A} / \mathfrak{m},\mathsf{A} / \mathfrak{m}) = 0\)（对 A 的每个极大理想 \(\mathfrak{m}\)）；
\item
  对 A 的每个极大理想 m，有 dp \(_{A}\) (A/m) \textless{} n。
\end{enumerate}

特别地，我们有 \(\mathrm{dh}(\Lambda)=\sup_{\mathfrak{m}}\mathrm{dp}_{\mathrm{A}}(\Lambda/\mathfrak{m})\)，其中 m 取遍 A 的极大理想组成的集合。

命题 6. 设 A 为 Noether 环，N 为 A-模，n 为一个 ≥ 0 的整数。下列条件等价：

\begin{enumerate}
\def\labelenumi{(\roman{enumi})}
\item
  有 \(\mathrm{di}_{\mathrm{A}}(\mathrm{N}) < n\)；
\item
  对每个素理想 \(\mathfrak{p}\)（属于 \(A\)），有 \(\mathrm{Ext}_{\mathrm{A}}^{n}(A/\mathfrak{p},\mathrm{N})=0\)；
\item
  对 A 的每个素理想 \(\mathfrak{p}\)，有 \(\operatorname{Ext}_{\mathrm{A}_{\mathfrak{p}}}^{n}(\kappa(\mathfrak{p}),\mathrm{N}_{\mathfrak{p}})=0\)。
\end{enumerate}

此外，若 A-模 N 有限生成，则这些条件等价于：

\begin{enumerate}
\def\labelenumi{(\roman{enumi})}
\setcounter{enumi}{3}
\tightlist
\item
  对 A 的每个极大理想 m，有 \(\operatorname{Ext}_{\mathrm{A}}^{i}(\mathrm{A}/\mathfrak{m},\mathrm{N})=0\) 对 \(n\leqslant i\leqslant n+\operatorname{ht}(\mathfrak{m})\) 成立。
\end{enumerate}

注意到条件 (iii) 等价于

(iii') 对 A 的每个素理想 \(\mathfrak{p}\)，有 \(\operatorname{Ext}_{\mathrm{A}}^{n}(\mathrm{A}/\mathfrak{p},\mathrm{N})\otimes_{\mathrm{A}}\kappa(\mathfrak{p})=0\)。

事实上，由于 \(\operatorname{Ext}_{\mathrm{A}}^{n}(\mathrm{A}/\mathfrak{p},\mathrm{N})\) 被 \(\mathfrak{p}\) 零化，A-模 \(\operatorname{Ext}_{\mathrm{A}}^{n}(\mathrm{A}/\mathfrak{p},\mathrm{N})\otimes_{\mathrm{A}}\kappa(\mathfrak{p})\) 同构于 \(\operatorname{Ext}_{\mathrm{A}}^{n}(\mathrm{A}/\mathfrak{p},\mathrm{N})\otimes_{\mathrm{A}}\mathrm{A}_{\mathfrak{p}}\)，从而同构于 \(\operatorname{Ext}_{\mathrm{A}_{\mathfrak{p}}}^{n}(\kappa(\mathfrak{p}),\mathrm{N}_{\mathfrak{p}})\) ( \(n^{\circ}\) 2, 命题 2)。

蕴含关系 (i)⇒(ii) 与 (i)⇒(iv) 由第 1 节的命题 1 得出，而蕴含关系 (ii)⇒(iii') 是显然的。

\begin{enumerate}
\def\labelenumi{(\roman{enumi})}
\setcounter{enumi}{2}
\tightlist
\item
  (iii)：对每个 A-模 M，令 T(M) = Ext \(_{A}^{n}\) (M,N)。设 (i) 不成立。则（第 1 节, 命题 1）存在 A 的理想 a 使 T(A/a) ≠ 0。设 p 为具有这一性质的理想中的极大者；我们来证明 p 是素理想。由 IV, § 1, 第 4 节, 定理 1 与定理 2，存在 A/p 的一个合成列 (M \(_{i}\) ) \(_{0\leqslant i\leqslant m}\)，使得每个商 M \(_{i}\) /M \(_{i+1}\) 都同构于某个模 A/p \(_{i}\)，其中 p \(_{i}\) (0 ≤ i ≤ m - 1) 是包含 p 的素理想。若对一切 i 都有 T(A/p \(_{i}\) ) 为零，则由正合列
\end{enumerate}

\[
\mathrm{T} \left(\mathrm{M} _ {0} / \mathrm{M} _ {i}\right) \longrightarrow \mathrm{T} \left(\mathrm{M} _ {0} / \mathrm{M} _ {i + 1}\right) \longrightarrow \mathrm{T} \left(\mathrm{M} _ {i} / \mathrm{M} _ {i + 1}\right)
\]

即可对 i 用归纳法推出 \(\mathrm{T}(\mathrm{M}_{0}/\mathrm{M}_{m})=\mathrm{T}(\mathrm{A}/\mathfrak{p})\) 为零，而事实并非如此。因此存在一个指标 i 使 \(\mathrm{T}(\mathrm{A}/\mathfrak{p}_{i})\neq0\)。由 p 的极大性，有 \(p=p_{i}\)，故 p 是素理想。

设 \(x\) 为 \(\mathrm{A} - \mathfrak{p}\) 的一个元素；由正合列

\[
0 \rightarrow \mathrm{A} / \mathfrak {p} \stackrel {{x _ {\mathrm{A} / \mathfrak {p}}}} {{\longrightarrow}} \mathrm{A} / \mathfrak {p} \longrightarrow \mathrm{A} / (\mathfrak {p} + x \mathrm{A}) \rightarrow 0
\]

得到正合列

\[
\mathrm{T} (\mathrm{A} / (\mathfrak {p} + x \mathrm{A})) \longrightarrow \mathrm{T} (\mathrm{A} / \mathfrak {p}) \stackrel {{u}} {{\longrightarrow}} \mathrm{T} (\mathrm{A} / \mathfrak {p}),
\]

其中 \(u = \operatorname{Ext}_{\mathrm{A}}^{n}(x_{\mathrm{A}}, 1_{\mathrm{N}}) - x_{\mathrm{T(A/p)}}\) (A, X, 第 89 页, 命题 6)。由 \(\mathfrak{p}\) 的极大性，有 \(\mathrm{T(A / (\mathfrak{p} + xA))} = 0\)，故同态 \(u\) 是单射。于是非零 A/p-模 \(\mathrm{T(A/p)}\) 是无扭的。这蕴含 \(\mathrm{T(A/p)} \otimes_{\mathrm{A}} \kappa(\mathfrak{p})\) 非零 (A, II, 第 117 页, 推论 1)，与 (iii') 矛盾。

设 A-模 N 有限生成，我们来证明 (iv) 蕴含 (iii)。设 p 为 A 的一个素理想，m 为包含 p 的一个极大理想。存在 A 的一个以 p 与 m 为端点的饱和素理想链。该链的长度 r 小于 ht(m)；因此在假设 (iv) 下，我们有

\[
\operatorname{Ext} _ {\mathrm{A} _ {\mathfrak {m}}} ^ {n + r} (\kappa (\mathfrak {m}), \mathrm{N} _ {\mathfrak {m}}) = \operatorname{Ext} _ {\mathrm{A}} ^ {n + r} (\mathrm{A} / \mathfrak {m}, \mathrm{N}) \otimes_ {A} \mathrm{A} _ {\mathfrak {m}} = 0.
\]

于是由 § 1, 第 7 节的引理 3 得出 \(\mathrm{Ext}_{\mathrm{A}_{\mathfrak{p}}}^{n}(\kappa(\mathfrak{p}),\mathrm{N}_{\mathfrak{p}})=0\)，这就证明了 (iii)。

注 5.— 设 N 为有限生成 A-模；条件「\(\operatorname{Ext}_{\mathrm{A}}^{n}(A/\mathfrak{m},\mathrm{N})=0\)（对 A 的每个极大理想 \(\mathfrak{m}\)）」并不必然蕴含 \(\operatorname{di}_{\mathrm{A}}(\mathrm{N})<n\)。例如，若环 A 是局部环且不是 Gorenstein 环 (第 7 节, 定义 1)，则有 \(\operatorname{Ext}_{A}^{n}(A/\mathfrak{m},\mathrm{A})=0\)（对 \(n<\operatorname{prof}(\mathrm{A})\)），但 \(\operatorname{di}_{\mathrm{A}}(\mathrm{A})=+\infty\)。

